\section{Discussion}
\label{sec:discussion}

The principal implication is methodological. Forward-only evaluation can certify
a fine-tuned model that has quietly lost capability on the inverse task. In
software engineering that risk spans more than one benchmark. Developers move
between paired tasks such as translation and back-translation and compilation and
decompilation. Mutation and repair raise a similar operational risk.
They fall outside Definition~\ref{def:pair}, because a mutant changes program
behavior. A pipeline that checks only the trained direction can promote the model variant that is least reliable once the workflow demands reversal, inspection, or rollback.

The risk extends beyond obfuscation. Security workflows couple forward and
reverse transformations, and recent work delegates both decompilation and
obfuscation to language models~\cite{tan2024llm4decompile, mohseni2025obfuscate}.
Modernization workflows do the same for translation, where forward translations
already carry functional bugs~\cite{pan2024lost, eniser2024rust}. Our checks on
translation and compilation show the same directional loss in those domains. The
governing claim concerns validation discipline. When a deployment context
contains a meaningful inverse task, that inverse must be evaluated in its own
right.

The mechanism evidence warrants a narrow reading. On Identifier Renaming, the
chain-of-thought probe agrees with a reasoning-generation gap. The model surfaces
the reverse identifier mapping yet generates the obfuscated code unchanged. The
behavior fits an account in which standard fine-tuning strengthens a forward
generation prior more than it preserves the bidirectional relation between source
and transformed code. The lexical result does not establish this as the sole
mechanism, and we do not claim that it carries unchanged to the harder tasks.

The recovery results fix the limits of what can be concluded. BFT shows that the
reverse direction is not inherently unreachable on the lexical and structural
tasks. Explicit supervision restores much of $\mathrm{Acc}_{\text{rev}}$ on
Identifier Renaming and Dead Code Insertion. CFT shows that part of the lexical
loss is recoverable without reverse labels, which points to a representational
component. Literal Encryption bounds these interventions. Full reverse
supervision fails there, so the present evidence cannot reduce the phenomenon to
one cause. The defensible summary is that directional capability loss interacts
with task complexity, and that the interaction becomes severe once the inverse
demands algorithmic recovery.

BRF follows from this finding as an instrument rather than a theory of model
cognition. It is a low-cost diagnostic that asks whether adaptation to a forward
transformation has introduced a hidden regression in the reverse direction. The
logic is the logic of regression testing. A local improvement is not sufficient
evidence of system reliability when the same change can break another part of the
workflow, and bidirectional evaluation supplies the missing check.

Several threats to validity remain. The main experiments use Java and six
decoder-only models, which limits the external validity of the exact percentages.
We mitigate this by spanning open and proprietary models and by adding the
translation and compilation checks. The strongest evidence of genuine capability
loss appears on Identifier Renaming, where the GPT models show measurable
baseline transformation ability and lose the reverse direction after standard
fine-tuning. On the harder tasks, the data supports the behavioral claim of
directional collapse more firmly than the stronger claim that an
already-demonstrated reverse skill was destroyed.

The study also stops short of direct representation analysis. We do not measure
internal geometry, activation dynamics, or decoding-time competition between
forward and reverse continuations, and those analyses are the natural next step.
Evaluation on encoder-decoder models, retrieval-augmented systems, and hybrid
symbolic-neural pipelines is also needed. The present results justify that
effort, since fine-tuning can reshape a code model in a directionally asymmetric
way that ordinary evaluation leaves invisible.
